\chapter{Polynomial consequences and proof degree}\label{tx:ch:consequences}
A polynomial equation can be used in two different ways. One can evaluate
it at a proposed solution, or one can combine it algebraically with other
equations. This book studies the second activity. Given equations
$f=0$ for $f$ in a set $\Gamma$, how complicated is a derivation of a
specified equation $G=0$? Our measure of complexity is the largest total
degree needed in the derivation.

The distinction between an equation and its left-hand polynomial is
worth keeping in view. We usually write $f$ instead of $f=0$, and write
$\Gamma\vdash G$ for deriving $G=0$. Thus a derivation of the polynomial
$1$ is a contradiction, since it asserts $1=0$. A derivation of $x+y$
need not be a contradiction: it may simply express a consequence shared
by all solutions of the original equations.

\section{The ambient polynomial ring}
Let $K$ be a field and let $V$ be a finite set of variable names. We work
in the ordinary polynomial ring $K[x_v:v\in V]$. A monomial is a product
$\prod_vx_v^{a_v}$ with nonnegative integer exponents. Its total degree
is $\sum_va_v$. The total degree of a nonzero polynomial is the largest
degree of a monomial with nonzero coefficient, after equal monomials
have been collected. We use the convention $\deg0=0$.

For example, over $\mathbb F_2$ the polynomial
\[
                   (xy+x)+(xy+1)=x+1
\]
has degree one. The two degree-two terms have cancelled. In a polynomial
ring over a field, nonzero products satisfy
$\deg(fg)=\deg f+\deg g$. We will use this fact when we define the
degree of NS certificates and PC derivations in the next two sections.

A system $\Gamma$ consists of polynomials regarded as equations equal
to zero. It may include Boolean equations $x_v^2-x_v$, but their presence
does not change the ambient ring. In particular, $x_v^2$ is still a
polynomial of degree two. Replacing it by $x_v$ requires a justified
algebraic step using $x_v^2-x_v$.

\section{Nullstellensatz certificates}
The algebraic proof system was introduced by Beame, Impagliazzo,
Kraj\'i\v cek, Pitassi and Pudl\'ak~\cite{NSFoundations}. We use its
polynomial-identity form, with degree measured on the complete products.
\begin{definition}[NS consequence]\label{tx:def:NS}
An NS derivation of $G$ from $\Gamma$ is a finite polynomial identity
\begin{equation}\label{tx:eq:NS}
                         G=\sum_{f\in F}a_ff,
                   \qquad F\subseteq\Gamma\text{ finite}.
\end{equation}
Its degree is the largest degree of a nonzero product $a_ff$.
Thus it has degree at most $d$ if
$\deg a_f+\deg f\le d$ for every nonzero summand. Zero summands can be
removed. The empty identity is a degree-zero derivation of $G=0$.
\end{definition}

The letters NS refer to Nullstellensatz, but the definition is simply
an identity which can be checked by multiplication and addition. There
is no appeal to an existence theorem hidden in a certificate. The
multiplier $a_f$ records how the input equation $f=0$ is used.

\begin{example}[A first certificate]\label{tx:eg:first-NS}
From $f_1=x+y$ and $f_2=y-z$, over any field, we derive $G=x+z$ by
\[
                       x+z=f_1-f_2.
\]
Both products have degree one. From $f_1=x$ and $f_2=1-x$ we instead
obtain $1=f_1+f_2$. The first identity expresses a consequence; the
second is a refutation.
\end{example}

The relevant degree is the degree of each product before the products
are added. A large-degree contribution cannot be ignored merely because
it cancels against another summand in the final identity. Otherwise
every refutation would appear to have degree zero, since its final
polynomial is the constant $1$.

\section{Ordinary Polynomial Calculus}
Polynomial Calculus was introduced by Clegg, Edmonds and
Impagliazzo~\cite{PCFoundations}. The rules below fix the ordinary
polynomial convention used throughout this book.
\begin{definition}[Ordinary-PC consequence]\label{tx:def:PC}
A PC derivation of $G$ is a finite sequence $P_1,\ldots,P_M$ ending in
$G$. A line may be an input polynomial from $\Gamma$, a linear
combination $\alpha P_i+\beta P_j$ of earlier lines, or a variable
multiple $x_vP_i$ of an earlier line. Coefficients $\alpha,\beta$ belong
to $K$. The zero polynomial is allowed as the empty linear combination.
The degree of the derivation is $\max_i\deg P_i$, with each line
collected as an actual polynomial.
\end{definition}

The word \emph{ordinary} stresses that there is no automatic Boolean
reduction and no rule which multiplies by an arbitrary polynomial in a
single step. A monomial multiple can be obtained by repeated variable
multiplication, but every intermediate line counts towards the degree.
The number of lines is unrestricted.

\begin{example}[A cancellation which enables the next step]\label{tx:eg:cancellation}
Over $\mathbb F_2$, consider the inputs
\[
                   f_1=xy+x,\qquad f_2=xy+1,\qquad f_3=xz.
\]
We want a PC derivation of $z$ in which every line has degree at most
two. This is a restriction on the whole proof, including its inputs and
intermediate lines. After listing these inputs, PC may continue with
\begin{align*}
 f_1+f_2&=x+1,\\
 z(x+1)&=xz+z,\\
 (xz+z)+f_3&=z.
\end{align*}
Every line stays within the chosen degree-two budget. Multiplication
by $z$ is always a PC rule, but here we must also check that its result
fits that budget. The preceding line $x+1$ has degree one, so
$z(x+1)=xz+z$ has degree two and is an admissible line in this
degree-two derivation. Multiplying one of the original quadratic inputs
by $z$ would instead produce a degree-three line. Expanding the same calculation into
one NS identity gives
\[
                          z=zf_1+zf_2+f_3,
\]
whose first two products have degree three. This is a concrete reason to
keep track of cancellations inside a PC proof.

The axioms are consistent: $x=y=1$ and $z=0$ is a solution. We have
proved the nonconstant consequence $z=0$, not a contradiction.
Exercise~\ref{tx:ex:small-gap} checks that the two minimum degrees in
this example really are different.
\end{example}

More generally, suppose we require every line to have degree at most
$d$, and the proof has derived a nonzero polynomial $P$. A variable
multiple stays within this degree budget only if $\deg P\le d-1$. This follows from
$\deg(x_vP)=\deg P+1$. Even if $P$ is represented as a linear combination
of large-degree earlier lines, it is its actual degree \emph{after
cancellation} that controls this next step. This observation is the
source of the intersection operation in our later linear-algebra model.

\section{Least degrees, zero, and infinity}
\begin{definition}[Least target degree]\label{tx:def:least-degree}
For $S\in\{\NS,\PC\}$, let $h_S(\Gamma\vdash G)$ be the least degree
of an $S$ derivation of $G$, or $\infty$ if there is no derivation.
We order $\NN\cup\{\infty\}$ by placing $\infty$ above every integer.
For refutations write
\[
                      d_S(\Gamma)=h_S(\Gamma\vdash1).
\]
\end{definition}

Thus $h_{\PC}(\Gamma\vdash G)=2$ says both that a PC derivation with
every line of degree at most two exists and that no derivation confined
to degree at most one exists. The symbol $S$ selects the proof system;
$h_S$ then records the smallest degree achievable in that system.

These conventions separate three situations. The zero target has degree
zero, even from empty input. A nonzero constant target is equivalent to
the target $1$, by rescaling with a field element. A nonconstant target
has a positive degree floor, but its minimum proof degree can be much
larger or infinite.

\begin{proposition}[Degree floor and simulation]\label{tx:prop:basic-comparison}
For every target $G$,
\begin{equation}\label{tx:eq:basic-comparison}
 \deg G\le h_{\PC}(\Gamma\vdash G)
          \le h_{\NS}(\Gamma\vdash G).
\end{equation}
If a finite proof exists, the first inequality is an ordinary integer
inequality. With an infinite value it is read in the extended order.
\end{proposition}
\begin{proof}
The last PC line is $G$, so its degree is at least $\deg G$. For NS,
a sum of products of degree at most $d$ is also of degree at most $d$.
To simulate an NS certificate, expand each multiplier into monomials.
Derive a permitted product $\mu f$ by listing $f$ and multiplying by
one variable factor of $\mu$ at a time. None of these intermediate
products exceeds $\deg\mu+\deg f$. Rescale and add all products to
obtain $G$. Every step remains within the original NS degree bound.
\end{proof}

The algebraic object which organises unrestricted derivability is an
\emph{ideal}. An ideal in a polynomial ring is a subspace closed under
multiplication by every polynomial. The ideal generated by $\Gamma$,
written $\langle\Gamma\rangle$, consists of all finite sums
$\sum_f a_ff$ with $f\in\Gamma$ and polynomial coefficients $a_f$.
It is the smallest ideal containing the inputs. Thus unrestricted NS
derivability is precisely membership in $\langle\Gamma\rangle$.
At a fixed degree, the restriction on each product is an additional
condition; an arbitrary ideal representation need not respect it.

There is a converse without a degree bound. Annotate each PC input $f$
with its expression as a single input multiple. Linear combination of
two lines combines their annotations, and multiplication by $x_v$
multiplies every coefficient in the annotation by $x_v$. At the last
line this produces an NS identity. It may have a higher product degree
than the largest actual PC line, because cancellations have been
expanded rather than reused. Consequently NS and PC derive precisely
the same polynomials when degree is unrestricted: the ideal generated
by $\Gamma$. The difficult comparison later in the book concerns their
\emph{bounded-degree} behaviour uniformly in the universe size.

\section{Algebraic consequence and solutions}
Every derivation is sound under evaluation. If a tuple $a$ satisfies
$f(a)=0$ for every input, an NS identity gives $G(a)=0$. The PC rules
preserve this assertion line by line. The same is true when $a$ lies
in an extension field of $K$.

The converse need not hold for the proof systems defined here. From
$x^2=0$, every field-valued solution has $x=0$, but $x$ is not a
polynomial multiple of $x^2$. It therefore has no NS or PC derivation.
So there are two different questions: does $G$ vanish at every solution,
and does an NS or PC derivation of $G$ exist? We call the first
\emph{semantic implication} and the second \emph{algebraic derivability}.
For the proof systems defined here, algebraic derivability is ideal
membership. The example shows that semantic implication alone does not
guarantee it.

For Boolean systems we can prove a converse for refutation. Here a
\emph{Boolean assignment} assigns each variable either zero or one, and
an \emph{inconsistent Boolean system} means that no such assignment
satisfies all its equations. We now prove completeness: an inconsistent
Boolean system with the Boolean axioms explicitly present has an
algebraic refutation. The argument constructs the certificate directly.

\begin{proposition}[Completeness for Boolean refutation]\label{tx:prop:boolean-completeness}
Let $\Gamma$ be a finite set of polynomials in finitely many variables,
over any field $K$, containing $x_v^2-x_v$ for every variable $x_v$.
If no Boolean assignment satisfies $\Gamma$, there is an NS identity
deriving $1$ from $\Gamma$, and hence a PC refutation.
\end{proposition}
\begin{proof}
The Boolean axiom gives $x_v(x_v-1)=0$ at every solution. Since a field
has no zero divisors, each coordinate is zero or one; this remains true
over an extension field. For each Boolean assignment $a$, put
\[
           \delta_a=\prod_{a_v=1}x_v\prod_{a_v=0}(1-x_v).
\]
This is its indicator polynomial: $\delta_a(a)=1$, while
$\delta_a(b)=0$ for every other Boolean assignment $b$. Expanding
$\prod_v(x_v+(1-x_v))$ gives $\sum_a\delta_a=1$.

There are only $2^{|V|}$ Boolean assignments. Because none satisfies the
system, for each $a$ choose one input $f_a$ with $f_a(a)\ne0$. Its
nonzero value has an inverse in $K$. Form the explicit polynomial
\[
                  H=1-\sum_a\frac{\delta_a}{f_a(a)}f_a.
\]
At any Boolean assignment $b$, all summands except the one indexed by
$b$ vanish, and that summand has value one. Thus $H(b)=0$.

We claim that a polynomial vanishing on every Boolean assignment is
an explicit sum of multiples of the Boolean axioms. To see this,
repeatedly reduce powers using the identity
\[
 x_v^e-x_v=(x_v^2-x_v)(1+x_v+\cdots+x_v^{e-2})\qquad(e\ge2).
\]
Each replacement records a polynomial multiple of $x_v^2-x_v$.
After finitely many replacements every exponent is at most one, so
\[
                  H=R+\sum_v b_v(x)(x_v^2-x_v),
\]
where $R$ is multilinear. The replacements do not change values on
Boolean assignments, so $R$ also vanishes at all of them. A multilinear
polynomial is determined by these values. For a direct induction, write
$R=A+x_v B$: evaluating at $x_v=0$ gives that $A$ vanishes on the
smaller Boolean cube, and evaluating at $x_v=1$ then gives the same for
$B$. Starting with constant polynomials, induction proves $R=0$.
Therefore rearranging the definition of $H$ gives the actual identity
\[
       1=\sum_a\frac{\delta_a}{f_a(a)}f_a
                            +\sum_v b_v(x)(x_v^2-x_v).
\]
Every polynomial on the right which is used as an input belongs to
$\Gamma$. This is an NS refutation. The simulation proved earlier
converts it into a PC refutation.
\end{proof}

Completeness here means that absence of a satisfying assignment
guarantees existence of a refutation. The construction uses an
exhaustive list of assignments and power reductions, so it gives an
ordinary proof with finitely many terms or lines. It does not promise
that the proof is short or has low degree. The degree lower bounds
later in the book concern how much degree may be unavoidable despite
this existence result.

\section{The next question: can a proof survive a change of universe?}
The definitions so far concern one fixed polynomial system. In the
next chapter the axioms and targets will come from fixed rules indexed
by the universe size $n$. We shall then ask whether a degree-$d$ proof
at one size can be lifted to a larger size with the same degree budget.
The earlier proof may use many lines; our aim is to preserve the
degree bound without imposing a bound on their number.

If this can be repeated at $n,n+q,n+2q,\ldots$, the one chosen proof
degree becomes a constant bound on that progression. If the family
is known to have no such bounded progression, an eligible starting
proof cannot exist. Measuring which degrees are eligible is how the
lifting construction will become a lower-bound argument.

\section{Exercises}
\begin{exercise}\label{tx:ex:zero-constant}
Find the least NS and PC degrees of the targets $0$, $1$, and $x$ from
$\Gamma=\varnothing$ and from $\Gamma=\{1\}$. Explain the role of the
last-line degree floor.
\end{exercise}
\noindent{\small\hyperref[tx:ans:zero-constant]{Model answer (p.~\pageref*{tx:ans:zero-constant}).}}\par\medskip

\begin{exercise}\label{tx:ex:small-gap}
For Example~\ref{tx:eg:cancellation}, prove
$h_{\PC}(\Gamma\vdash z)=2$ and
$h_{\NS}(\Gamma\vdash z)=3$. For the NS lower bound, list the monomials
which can occur in a degree-two certificate.
\end{exercise}
\noindent{\small\hyperref[tx:ans:small-gap]{Model answer (p.~\pageref*{tx:ans:small-gap}).}}\par\medskip

\begin{exercise}\label{tx:ex:Boolean-rule}
A proposed degree-two PC calculation multiplies the derived line $xy$
by $x$ and replaces $x^2y$ by $xy$, using a Boolean equation. Explain
why this is not an ordinary degree-two step, even if $x^2-x$ is an input.
\end{exercise}
\noindent{\small\hyperref[tx:ans:Boolean-rule]{Model answer (p.~\pageref*{tx:ans:Boolean-rule}).}}\par\medskip

